\documentclass[11pt,a4paper]{article}
\usepackage[margin=1in]{geometry}
\usepackage{amsmath,amssymb,amsfonts,amsthm}
\usepackage{bm}
\usepackage{booktabs}
\usepackage{enumitem}
\usepackage{hyperref}
\usepackage{cleveref}
\usepackage{algorithm}
\usepackage{algpseudocode}

% ── Notation macros (consistent with Pythonic-DISORT.ipynb) ──────────────
\newcommand{\R}{\mathbb{R}}
\newcommand{\tr}{\operatorname{tr}}
\newcommand{\diag}{\operatorname{diag}}
\newcommand{\cond}{\operatorname{cond}}
\newcommand{\expm}{\operatorname{expm}}
\newcommand{\vect}[1]{\bm{#1}}
\newcommand{\mat}[1]{\mathbf{#1}}
\newcommand{\tbot}{\tau_{\mathrm{bot}}}
\newcommand{\dd}{\,\mathrm{d}}

\theoremstyle{definition}
\newtheorem{proposition}{Proposition}
\newtheorem{remark}{Remark}
\newtheorem{definition}{Definition}

\title{The Star-Product Magnus Integrator for Radiative Transfer\\
with Continuously $\tau$-Varying Optical Properties}
\author{Technical Report --- Extension to the Pythonic-DISORT\\
Comprehensive Documentation}
\date{March 2026}

\begin{document}
\maketitle

\begin{abstract}
This report documents the \texttt{pydisort\_magnus} solver, a forward
solver for the one-dimensional radiative transfer equation (RTE) in a
plane-parallel atmosphere with continuously $\tau$-varying
single-scattering albedo $\omega(\tau)$ and phase function.
It combines first-order Magnus integration (midpoint-rule matrix
exponential) with Redheffer star-product accumulation of
$N\!\times\!N$ reflection, transmission, and source operators.
The resulting algorithm is unconditionally numerically stable for
arbitrary optical depth, achieves $O(h^2)$ global convergence, and
requires no eigendecomposition.  We compare the Magnus approach
mathematically and computationally with multi-layer DISORT (as used
in SCIATRAN) and adding-doubling, discuss the fourth-order Magnus
extension and the principle of solver-function smoothness matching,
and outline an adaptive step-control strategy based on the matrix
commutator.  This report extends sections~3--4 of the Pythonic-DISORT
Comprehensive Documentation~\cite{HP2024doc} and assumes thorough
familiarity with that document.
\end{abstract}

\tableofcontents
\vspace{1em}

%======================================================================
\section{Introduction and Motivation}\label{sec:intro}
%======================================================================

The standard \texttt{pydisort} solver (sections~3--4 of the
Comprehensive Documentation) handles multi-layer atmospheres by
assuming piecewise-constant optical properties within each layer and
performing an eigendecomposition of the coefficient matrix
$\mat{A}$ in each layer.  This approach is exact for homogeneous
layers and numerically stable via the Stamnes--Conklin
substitution~\cite{stamnes1984}.

However, for atmospheres where the single-scattering albedo
$\omega(\tau)$ and the phase function $p(\nu;\tau)$ vary
\emph{continuously} with optical depth---as arises in cloud
microphysics when the effective radius $r_e$ varies with $\tau$---the
piecewise-constant approximation introduces $O(h^2)$ discretisation
error, where $h$ is the layer thickness.  While this error can be
reduced by increasing the number of layers, the cost of an
eigendecomposition per layer ($O(N^3)$ per layer per Fourier mode)
makes this approach expensive for fine discretisations.

The Magnus integrator~\cite{magnus1954,blanes2009} offers a
natural alternative: it directly handles the variable-coefficient
ODE $\vect{I}'(\tau)=\mat{A}(\tau)\vect{I}(\tau)+\vect{S}(\tau)$ via
a midpoint-rule matrix exponential at each step, without
eigendecomposition.  The first-order Magnus method is also $O(h^2)$
accurate, but each step costs only a single matrix exponential
(\texttt{scipy.linalg.expm}), which is cheaper than a full
eigendecomposition for the same matrix size.

The challenge is combining the per-step results across $K$ steps to
obtain the solution over the full optical depth $[0,\tbot]$.
Na\"ive accumulation of the $2N\!\times\!2N$ propagator fails for
thick atmospheres because $\cond(\mat{\Phi})=
e^{2\lambda_{\max}\tbot}$, which exceeds double precision for
$\tbot\gtrsim\ln(10^{16})/\lambda_{\max}\approx 2.6$ when
$\lambda_{\max}\approx 14$ (for $N\!=\!4$ streams).
\Cref{sec:propagator_problem} analyses this failure in detail.

The Redheffer star product~\cite{redheffer1962,grant1969}
resolves this by reformulating the problem in terms of
$N\!\times\!N$ reflection, transmission, and source operators,
all of which remain $O(1)$ regardless of optical depth.
\Cref{sec:star_product,sec:extraction,sec:algorithm} develop the
full algorithm, and \Cref{sec:bc_solver} derives the boundary
condition solver.  \Cref{sec:convergence} presents convergence
analysis and numerical verification.
\Cref{sec:comparison} places the algorithm in context by comparing
it with multi-layer DISORT/SCIATRAN and adding-doubling, both
mathematically and computationally.  \Cref{sec:higher_order} develops
the fourth-order Magnus extension and the principle of solver-function
smoothness matching, with application to adiabatic cloud profiles and
profiles perturbed by precipitation.  \Cref{sec:adaptivity} outlines
a commutator-based adaptive step-control strategy.


%======================================================================
\section{The Variable-Coefficient Discrete-Ordinate ODE}\label{sec:ode}
%======================================================================

We begin from the azimuthally decomposed, angularly discretised RTE
as derived in sections~3.2--3.4 of the Comprehensive Documentation,
but now with $\tau$-dependent optical properties.

\subsection{Fourier decomposition}

The diffuse intensity $u(\tau,\mu,\phi)$ is expanded in a Fourier
cosine series (cf.\ section~3.2 of the Documentation):
\begin{equation}\label{eq:fourier}
  u(\tau,\mu,\phi)
  = \sum_{m=0}^{M_F-1} u^m(\tau,\mu)\,\cos\bigl(m(\phi_0-\phi)\bigr),
\end{equation}
where $M_F$ is the number of Fourier modes.  Each mode $m$ is solved
independently.

\subsection{The coefficient matrix \texorpdfstring{$\mat{A}^m(\tau)$}{A\_m(tau)}}

After double-Gauss quadrature discretisation with $2N$ streams
(section~3.4 of the Documentation), each Fourier mode satisfies the
$2N$-dimensional linear ODE
\begin{equation}\label{eq:ode}
  \frac{d}{d\tau}\begin{pmatrix}\vect{u}^+\\\vect{u}^-\end{pmatrix}
  = \mat{A}^m(\tau)
    \begin{pmatrix}\vect{u}^+\\\vect{u}^-\end{pmatrix}
  + \vect{S}^m(\tau),
  \qquad \tau\in[0,\tbot],
\end{equation}
where $\vect{u}^+ = (u^m(\tau,\mu_1),\dots,u^m(\tau,\mu_N))^T$ collects the
upward intensities and $\vect{u}^- = (u^m(\tau,-\mu_1),\dots,u^m(\tau,-\mu_N))^T$
the downward intensities.  The coefficient matrix has the block structure
\begin{equation}\label{eq:A_block}
  \mat{A}^m(\tau) = \begin{pmatrix}
    -\mat{\alpha}^m(\tau) & -\mat{\beta}^m(\tau) \\
     \mat{\beta}^m(\tau)  &  \mat{\alpha}^m(\tau)
  \end{pmatrix},
\end{equation}
with
\begin{equation}\label{eq:alpha_beta}
  \mat{\alpha}^m(\tau)
  = \mat{M}^{-1}\bigl(\mat{D}^{m,+}(\tau)\,\mat{W} - \mat{I}_N\bigr),
  \qquad
  \mat{\beta}^m(\tau)
  = \mat{M}^{-1}\,\mat{D}^{m,-}(\tau)\,\mat{W},
\end{equation}
where $\mat{M}=\diag(\mu_1,\dots,\mu_N)$ and
$\mat{W}=\diag(w_1,\dots,w_N)$ are the quadrature cosine and weight
matrices, and the phase-function kernel matrices are
\begin{equation}\label{eq:D_pm}
  D^{m,+}_{ij}(\tau) = \omega(\tau)\,\mathcal{D}^m(\tau;\,\mu_i,\,\mu_j),
  \qquad
  D^{m,-}_{ij}(\tau) = \omega(\tau)\,\mathcal{D}^m(\tau;\,\mu_i,\,-\mu_j).
\end{equation}
Here $\mathcal{D}^m$ is the $\omega$-free phase-function kernel:
\begin{equation}\label{eq:D_pure}
  \mathcal{D}^m(\tau;\,\mu_i,\,\mu_j)
  = \frac{1}{2}\sum_{\ell=m}^{L-1}(2\ell+1)\,
    \frac{(\ell-m)!}{(\ell+m)!}\,
    g_\ell(\tau)\,P_\ell^m(\mu_i)\,P_\ell^m(\mu_j),
\end{equation}
where $g_\ell(\tau)$ are the (unweighted) Legendre coefficients of the
phase function at optical depth $\tau$, $P_\ell^m$ are the associated
Legendre functions, and $L$ is the truncation order.

\begin{remark}[Separation of $\omega$ from the kernel]
The factorisation $D^{m,\pm}_{ij}=\omega(\tau)\cdot\mathcal{D}^m$
in~\eqref{eq:D_pm} is a deliberate design choice: the user supplies
$\omega(\tau)$ as a scalar function and $\mathcal{D}^m$ as a separate
callable, enabling independent variation of absorption and scattering
anisotropy.
\end{remark}

\subsection{The source vector \texorpdfstring{$\vect{S}^m(\tau)$}{S\_m(tau)}}

For a collimated beam of intensity $I_0$ incident at cosine angle
$\mu_0$ and azimuth $\phi_0$, the source vector for Fourier mode $m$
is (cf.\ section~3.6.1 of the Documentation):
\begin{equation}\label{eq:source}
  \vect{S}^m(\tau) = \begin{pmatrix} -\vect{Q}^{m,+}(\tau)\\
  \vect{Q}^{m,-}(\tau) \end{pmatrix},
\end{equation}
where the $i$-th component of $\vect{Q}^{m,+}$ is
\begin{equation}\label{eq:Q_plus}
  Q^{m,+}_i(\tau)
  = \frac{I_0}{4\pi}\,(2-\delta_{m0})\,
    \frac{\omega(\tau)}{\mu_i}\,
    \mathcal{D}^m(\tau;\,\mu_i,\,-\mu_0)\,
    e^{-\tau/\mu_0},
\end{equation}
and $Q^{m,-}_i(\tau)$ is the same expression with $\mu_i$ replaced by
$-\mu_i$ in the $\mathcal{D}^m$ evaluation (and sign
reversal as in section~3.4 of the Documentation).


%======================================================================
\section{First-Order Magnus Integration}\label{sec:magnus}
%======================================================================

\subsection{The Magnus expansion}

The Magnus expansion~\cite{magnus1954} writes the solution of the
matrix ODE $\vect{I}'=\mat{A}(\tau)\vect{I}$ as
$\vect{I}(\tau)=\exp\bigl(\mat{\Omega}(\tau)\bigr)\vect{I}(0)$, where
\begin{equation}
  \mat{\Omega}(\tau) = \mat{\Omega}_1(\tau)
  + \mat{\Omega}_2(\tau) + \cdots
\end{equation}
is an infinite series in nested integrals and commutators of $\mat{A}$.
The first-order term is $\mat{\Omega}_1(\tau)=\int_0^\tau\mat{A}(s)\dd s$.
Truncation at $\mat{\Omega}_1$ and approximation by the midpoint
quadrature rule gives the first-order Magnus method:
\begin{equation}\label{eq:magnus_step}
  \mat{\Phi}_k = \exp\!\bigl(h\,\mat{A}(\tau_{k+1/2})\bigr),
  \qquad \tau_{k+1/2}=(k+\tfrac{1}{2})h,
  \qquad h=\tbot/K,
\end{equation}
where $K$ is the number of equidistant steps.

The convergence radius of the Magnus series is
$\int_0^h\|\mat{A}(s)\|_2\dd s < \pi$~\cite{blanes2009,moan2008}.
For the DISORT coefficient matrix with
$\|\mat{A}\|_2\approx\lambda_{\max}\approx 1/\mu_{\min}$
(e.g.\ $\approx\!14.4$ for $N\!=\!4$), this requires
$h < \pi/\lambda_{\max}\approx 0.22$.  This bounds the
\emph{step size}, not the total depth: for any $\tbot$, choosing
$K > \lambda_{\max}\tbot/\pi$ suffices.

\subsection{The augmented matrix exponential}\label{sec:augmented}

To handle the source term $\vect{S}(\tau)$ in~\eqref{eq:ode}
simultaneously with the homogeneous propagator, we use the standard
augmented-system technique (cf.\ Van Loan~\cite{vanloan1978}).
Define the $(2N\!+\!1)\!\times\!(2N\!+\!1)$ matrix
\begin{equation}\label{eq:augmented}
  \mat{M}_k = \begin{pmatrix}
    h\,\mat{A}(\tau_{k+1/2}) & h\,\vect{S}(\tau_{k+1/2}) \\
    \vect{0}^T & 0
  \end{pmatrix}.
\end{equation}
Then
\begin{equation}\label{eq:augmented_expm}
  \exp(\mat{M}_k) = \begin{pmatrix}
    \mat{\Phi}_k & \vect{\delta}_k \\
    \vect{0}^T   & 1
  \end{pmatrix},
\end{equation}
where $\mat{\Phi}_k\in\R^{2N\times 2N}$ is the homogeneous step
propagator and $\vect{\delta}_k\in\R^{2N}$ is the particular-solution
increment from the source term over one step.  Both are obtained from
a single call to \texttt{scipy.linalg.expm}.

\subsection{Per-step accuracy}

The local truncation error of the first-order Magnus step is $O(h^3)$
(from the midpoint rule applied to $\mat{\Omega}_1$, plus the neglected
$\mat{\Omega}_2$ term which is $O(h^3)$).  Over $K=\tbot/h$ steps,
the global error is $O(h^2)$---i.e.\ second-order convergence.
This is confirmed experimentally in \Cref{sec:convergence}.


%======================================================================
\section{The Propagator Accumulation Problem}\label{sec:propagator_problem}
%======================================================================

\subsection{Eigenvalue structure of $\mat{A}$}

The block symmetry of~\eqref{eq:A_block} guarantees that the
eigenvalues of $\mat{A}$ come in exact $\pm\lambda$ pairs
(cf.\ section~3.5 of the Documentation):
\begin{equation}
  \lambda_1 > \lambda_2 > \cdots > \lambda_N > 0
  > -\lambda_N > \cdots > -\lambda_1.
\end{equation}
For isotropic scattering with $N\!=\!4$ and $\omega=0.5$:
$\lambda_{\max}\approx 13.77$, $\lambda_{\min}\approx 0.97$.

\subsection{Condition number of the propagator}

Na\"ive accumulation builds the full propagator
$\mat{\Phi}(0\to\tbot)=\prod_{k=0}^{K-1}\mat{\Phi}_k$, which for
constant $\mat{A}$ equals $\exp(\mat{A}\tbot)$.  Its condition number is
\begin{equation}\label{eq:cond_propagator}
  \cond\bigl(\exp(\mat{A}\tbot)\bigr)
  \geq \frac{\sigma_{\max}}{\sigma_{\min}}
  \geq \frac{e^{\lambda_{\max}\tbot}}{e^{-\lambda_{\max}\tbot}}
  = e^{2\lambda_{\max}\tbot}.
\end{equation}
For $\lambda_{\max}=14$ and $\tbot=5$:
$\cond\geq e^{140}\approx 10^{60.8}$, catastrophically exceeding
double precision ($\epsilon_{\text{mach}}\approx 10^{-16}$).

\subsection{Non-normality and transient growth}\label{sec:nonnormal}

The $\mat{M}^{-1}$ factor in~\eqref{eq:alpha_beta} makes $\mat{A}$
non-normal: $\mat{A}\mat{A}^T\neq\mat{A}^T\mat{A}$.
For a normal matrix, the singular values of $e^{\mat{A}\tau}$ would be
$e^{\text{Re}(\lambda_j)\tau}$---those with $\lambda_j>0$ grow while
those with $\lambda_j<0$ decay.  For the non-normal DISORT matrix,
however, non-orthogonality of the eigenvectors causes all singular
values to eventually grow as $\sigma_j\sim C_j e^{\lambda_{\max}\tau}$
for large $\tau$~\cite{kreiss1962}.  This
``dip-and-rebound'' phenomenon means that all $2N$ modes eventually
become large, precluding truncation of ``decaying'' modes.

\subsection{Why this is fundamental}

\emph{Any} algorithm that explicitly forms the $2N\!\times\!2N$
propagator $\mat{\Phi}(0\to\tbot)$ will encounter the same problem.
The solution is to avoid the propagator entirely and work with
$N\!\times\!N$ scattering operators that remain bounded.


%======================================================================
\section{The Interaction Principle and Redheffer Star Product}
\label{sec:star_product}
%======================================================================

\subsection{Scattering-matrix representation}

A slab occupying $[0,\tbot]$ is characterised by how it transforms
radiation incident from both sides.  Following Grant and
Hunt~\cite{grant1969} and the matrix-operator method of Plass et
al.~\cite{plass1973}, we define six $N\!\times\!N$ operators:
\begin{definition}[Slab scattering operators]\label{def:slab_ops}
Given a slab with reflection, transmission, and internal source
properties, the \emph{scattering-matrix representation} is the
system
\begin{align}
  \vect{u}^+(0) &= \mat{R}_\uparrow\,\vect{u}^-(0)
                  + \mat{T}_\uparrow\,\vect{u}^+(\tbot)
                  + \vect{s}_\uparrow, \label{eq:interaction_up}\\
  \vect{u}^-(\tbot) &= \mat{T}_\downarrow\,\vect{u}^-(0)
                      + \mat{R}_\downarrow\,\vect{u}^+(\tbot)
                      + \vect{s}_\downarrow, \label{eq:interaction_down}
\end{align}
where:
\begin{itemize}[nosep]
\item $\mat{R}_\uparrow\in\R^{N\times N}$:
  reflection of downwelling radiation, observed at the top;
\item $\mat{T}_\uparrow\in\R^{N\times N}$:
  upward transmission of upwelling radiation from the bottom to the top;
\item $\mat{T}_\downarrow\in\R^{N\times N}$:
  downward transmission of downwelling radiation from the top to the bottom;
\item $\mat{R}_\downarrow\in\R^{N\times N}$:
  reflection of upwelling radiation, observed at the bottom;
\item $\vect{s}_\uparrow\in\R^N$:
  upward source at the top from internal scattering of the
  collimated beam;
\item $\vect{s}_\downarrow\in\R^N$:
  downward source at the bottom from internal scattering of the
  collimated beam.
\end{itemize}
\end{definition}

\begin{remark}[Connection to the Documentation's notation]
In the piecewise-constant layer formulation of sections~3--4 of the
Documentation, the boundary-condition system is expressed in terms of
the eigenvectors $G^{\pm}$ and the Stamnes--Conklin-scaled
coefficients.  The scattering-matrix
operators~\eqref{eq:interaction_up}--\eqref{eq:interaction_down}
encode the same input--output relation in a basis-free form, without
requiring an eigendecomposition.
\end{remark}

\subsection{The Redheffer star product}

Consider two adjacent slabs: slab~1 (accumulated, on top, operators
with capital letters) and slab~2 (new step, below, operators with
lowercase).  The combined slab's operators are given by the
\emph{Redheffer star product}~\cite{redheffer1962}:

\begin{definition}[Redheffer star product]\label{def:star_product}
Let the resolvent be
\begin{equation}\label{eq:resolvent}
  \mat{E} = (\mat{I}_N - \mat{r}_\uparrow\,\mat{R}_\downarrow)^{-1}.
\end{equation}
Then the combined slab operators are:
\begin{align}
  \mat{R}_\uparrow^{\star}
    &= \mat{R}_\uparrow
     + \mat{T}_\uparrow\,\mat{E}\,\mat{r}_\uparrow\,\mat{T}_\downarrow,
  \label{eq:star_Rup}\\
  \mat{T}_\uparrow^{\star}
    &= \mat{T}_\uparrow\,\mat{E}\,\mat{t}_\uparrow,
  \label{eq:star_Tup}\\
  \mat{T}_\downarrow^{\star}
    &= \mat{t}_\downarrow\,(\mat{T}_\downarrow
     + \mat{R}_\downarrow\,\mat{E}\,\mat{r}_\uparrow\,\mat{T}_\downarrow),
  \label{eq:star_Tdn}\\
  \mat{R}_\downarrow^{\star}
    &= \mat{r}_\downarrow
     + \mat{t}_\downarrow\,\mat{R}_\downarrow\,\mat{E}\,\mat{t}_\uparrow,
  \label{eq:star_Rdn}\\
  \vect{s}_\uparrow^{\star}
    &= \vect{s}_\uparrow
     + \mat{T}_\uparrow\,\mat{E}\,
       (\mat{r}_\uparrow\,\vect{s}_\downarrow + \vect{s}_{\uparrow,k}),
  \label{eq:star_sup}\\
  \vect{s}_\downarrow^{\star}
    &= \vect{s}_{\downarrow,k}
     + \mat{t}_\downarrow\,
       (\vect{s}_\downarrow + \mat{R}_\downarrow\,\mat{E}\,
        (\mat{r}_\uparrow\,\vect{s}_\downarrow + \vect{s}_{\uparrow,k})).
  \label{eq:star_sdn}
\end{align}
\end{definition}

The physical interpretation is transparent:
$\mat{E}=(\mat{I}-\mat{r}_\uparrow\mat{R}_\downarrow)^{-1}$
is the resolvent for the multiple-reflection series between the
bottom face of slab~1 and the top face of slab~2.
Equation~\eqref{eq:star_Rup}, for example, says: the combined
upward reflection is the original reflection $\mat{R}_\uparrow$
plus radiation that transmits down through slab~1, reflects off
slab~2, bounces between the two slabs (resolvent $\mat{E}$), and
transmits back up through slab~1.

\subsection{Boundedness of intermediates}

\begin{proposition}\label{prop:bounded}
For a non-conservative atmosphere ($\omega(\tau)<1$ for all $\tau$):
\begin{enumerate}[nosep]
\item The reflection operators satisfy $\|\mat{R}_\uparrow\|_2<1$
  and $\|\mat{R}_\downarrow\|_2<1$ (by energy conservation---some
  energy is absorbed at each scattering).
\item The transmission operators decay:
  $\|\mat{T}\|_2\leq e^{-\lambda_N\tau}$ (governed by the
  smallest eigenvalue magnitude), bounded above by~1.
\item The resolvent~\eqref{eq:resolvent} is well-conditioned:
  $\cond(\mat{I}-\mat{r}_\uparrow\mat{R}_\downarrow)$ remains
  close to~1 because the product
  $\|\mat{r}_\uparrow\|\cdot\|\mat{R}_\downarrow\|<1$.
\item The source vectors $\vect{s}_\uparrow$, $\vect{s}_\downarrow$
  are $O(1)$ (bounded by the rescaled beam intensity).
\end{enumerate}
\end{proposition}

\textbf{Crucial consequence.}  All intermediate $N\!\times\!N$
quantities in the star product are $O(1)$.  No exponentially growing
or decaying numbers appear anywhere.  This is the fundamental
difference from the $2N\!\times\!2N$ propagator accumulation, and is
the reason the star-product formulation is unconditionally stable.

\subsection{Connection to invariant imbedding}

The scattering-matrix representation is closely related to the
\emph{invariant imbedding} formulation of Bellman, Kalaba, and
Prestrud~\cite{bellman1960}.  In that framework, the reflection
matrix $\mat{R}(\tau)$ for the slab $[0,\tau]$ satisfies the matrix
Riccati ODE
\begin{equation}\label{eq:riccati}
  \frac{d\mat{R}}{d\tau}
  = -\mat{\alpha}\mat{R} - \mat{R}\mat{\alpha}
    - \mat{R}\mat{\beta}\mat{R} - \mat{\beta}.
\end{equation}
The Riccati equation is the \emph{infinitesimal generator} of the
star-product semigroup: the exact Riccati step from $\tau$ to
$\tau+h$ is given by the M\"obius transformation
\begin{equation}\label{eq:mobius}
  \mat{R}(\tau+h)
  = (\mat{a}\,\mat{R}(\tau)+\mat{b})\,
    (\mat{c}\,\mat{R}(\tau)+\mat{d})^{-1},
\end{equation}
where $\bigl[\begin{smallmatrix}\mat{a}&\mat{b}\\
\mat{c}&\mat{d}\end{smallmatrix}\bigr]=\exp(h\mat{A})$.
The star product~\eqref{eq:star_Rup}--\eqref{eq:star_sdn}
generalises this to include transmission and source operators.


%======================================================================
\section{Extracting Slab Operators from a Magnus Step}
\label{sec:extraction}
%======================================================================

Each Magnus step produces a $2N\!\times\!2N$ homogeneous propagator
$\mat{\Phi}_k$ and a $2N$-vector particular increment
$\vect{\delta}_k$ via~\eqref{eq:augmented_expm}.  We now derive the
extraction of the six $N\!\times\!N$ scattering operators.

\subsection{Block decomposition}

Write
\begin{equation}\label{eq:block_decomp}
  \mat{\Phi}_k = \begin{pmatrix}\mat{a}&\mat{b}\\\mat{c}&\mat{d}\end{pmatrix},
  \qquad
  \vect{\delta}_k = \begin{pmatrix}\vect{\delta}^+\\\vect{\delta}^-\end{pmatrix},
\end{equation}
where $\mat{a},\mat{b},\mat{c},\mat{d}\in\R^{N\times N}$ and
$\vect{\delta}^+,\vect{\delta}^-\in\R^N$.  The propagator relates
the intensities at the step boundaries:
\begin{equation}
  \begin{pmatrix}\vect{u}^+(\tau_{k+1})\\\vect{u}^-(\tau_{k+1})\end{pmatrix}
  = \begin{pmatrix}\mat{a}&\mat{b}\\\mat{c}&\mat{d}\end{pmatrix}
    \begin{pmatrix}\vect{u}^+(\tau_k)\\\vect{u}^-(\tau_k)\end{pmatrix}
  + \begin{pmatrix}\vect{\delta}^+\\\vect{\delta}^-\end{pmatrix}.
\end{equation}

\subsection{Extraction formulas}

Solving for $\vect{u}^+(\tau_k)$ from the top row and substituting
into the interaction
form~\eqref{eq:interaction_up}--\eqref{eq:interaction_down}
(with the convention that $\tau_k$ is the top and $\tau_{k+1}$
the bottom of the step slab), we obtain:
\begin{align}
  \mat{t}_\uparrow   &= \mat{a}^{-1},
  &\mat{r}_\uparrow   &= -\mat{a}^{-1}\mat{b},
  \label{eq:extract_up}\\
  \mat{r}_\downarrow  &= \mat{c}\,\mat{a}^{-1},
  &\mat{t}_\downarrow  &= \mat{d}-\mat{c}\,\mat{a}^{-1}\mat{b},
  \label{eq:extract_down}\\
  \vect{s}_\uparrow   &= -\mat{a}^{-1}\vect{\delta}^+,
  &\vect{s}_\downarrow &= \vect{\delta}^-
                         - \mat{c}\,\mat{a}^{-1}\vect{\delta}^+.
  \label{eq:extract_source}
\end{align}

\begin{proposition}[Conditioning of $\mat{a}$]\label{prop:cond_a}
For a single Magnus step of size $h$:
\begin{equation}
  \cond(\mat{a}) \approx 1 + O(\lambda_{\max}h).
\end{equation}
Specifically, for the DISORT matrix with $\lambda_{\max}\approx 14$
and typical step size $h\leq 0.05$:
$\cond(\mat{a})\leq 1.30$.  The extraction~\eqref{eq:extract_up}--\eqref{eq:extract_source}
is therefore numerically stable.
\end{proposition}

\subsection{Efficient batched solve}\label{sec:batched_solve}

All six operators
in~\eqref{eq:extract_up}--\eqref{eq:extract_source} depend on
$\mat{a}^{-1}$ applied to different right-hand sides.  Rather than
computing $\mat{a}^{-1}$ explicitly, we solve the single
$N\!\times\!(2N\!+\!1)$ system
\begin{equation}\label{eq:batched_extract}
  \mat{a}\,\mat{X}
  = \bigl[\,\mat{I}_N \;\big|\; \mat{b} \;\big|\; \vect{\delta}^+\,\bigr],
\end{equation}
yielding
$\mat{X}=[\mat{a}^{-1}\mid\mat{a}^{-1}\mat{b}\mid\mat{a}^{-1}\vect{\delta}^+]$,
from which all operators are read off directly.  This costs a single
\texttt{numpy.linalg.solve} call (an $N\!\times\!N$ LU factorisation
with $2N\!+\!1$ right-hand sides), which is more efficient and
numerically stabler than explicit inversion.


%======================================================================
\section{The Complete Algorithm}\label{sec:algorithm}
%======================================================================

\subsection{Initialisation}

The identity slab (zero-thickness, no scattering, no source) has
operators
\begin{equation}\label{eq:identity_slab}
  \mat{R}_\uparrow = \mat{0},\quad
  \mat{T}_\uparrow = \mat{I}_N,\quad
  \mat{T}_\downarrow = \mat{I}_N,\quad
  \mat{R}_\downarrow = \mat{0},\quad
  \vect{s}_\uparrow = \vect{0},\quad
  \vect{s}_\downarrow = \vect{0}.
\end{equation}

\subsection{Accumulation loop}

\begin{algorithm}[H]
\caption{Star-product Magnus accumulation}\label{alg:accumulation}
\begin{algorithmic}[1]
\Require $\mat{A}^m(\cdot)$, $\vect{S}^m(\cdot)$, $\tbot$, $K$, $N$
\State Initialise $(\mat{R}_\uparrow,\mat{T}_\uparrow,\mat{T}_\downarrow,
       \mat{R}_\downarrow,\vect{s}_\uparrow,\vect{s}_\downarrow)$
       as in~\eqref{eq:identity_slab}
\State $h \gets \tbot/K$
\For{$k=0,\dots,K-1$}
  \State $\tau_{\text{mid}} \gets (k+\tfrac{1}{2})h$
  \State Build augmented matrix $\mat{M}_k$ per~\eqref{eq:augmented}
  \State $\exp(\mat{M}_k) \to (\mat{\Phi}_k,\,\vect{\delta}_k)$
         \Comment{Single \texttt{expm} call}
  \State Extract $(\mat{r}_\uparrow,\mat{t}_\uparrow,\mat{t}_\downarrow,
         \mat{r}_\downarrow,\vect{s}_{\uparrow,k},\vect{s}_{\downarrow,k})$
         via~\eqref{eq:batched_extract}
  \State Update accumulated operators via
         star product~\eqref{eq:star_Rup}--\eqref{eq:star_sdn}
\EndFor
\State \Return $(\mat{R}_\uparrow,\mat{T}_\uparrow,\mat{T}_\downarrow,
       \mat{R}_\downarrow,\vect{s}_\uparrow,\vect{s}_\downarrow)$
\end{algorithmic}
\end{algorithm}

\begin{remark}[Efficient resolvent computation]
The star-product update requires the resolvent
$\mat{E}=(\mat{I}-\mat{r}_\uparrow\mat{R}_\downarrow)^{-1}$
applied to multiple vectors.  As with the extraction step, we
form a single $N\!\times\!(2N\!+\!1)$ batched solve:
\begin{equation}\label{eq:batched_star}
  (\mat{I}-\mat{r}_\uparrow\mat{R}_\downarrow)\,\mat{Y}
  = \bigl[\,\mat{r}_\uparrow\mat{T}_\downarrow
    \;\big|\; \mat{t}_\uparrow
    \;\big|\; \mat{r}_\uparrow\vect{s}_\downarrow
              + \vect{s}_{\uparrow,k}\,\bigr],
\end{equation}
yielding $\mat{E}\mat{r}_\uparrow\mat{T}_\downarrow$,
$\mat{E}\mat{t}_\uparrow$, and
$\mat{E}(\mat{r}_\uparrow\vect{s}_\downarrow+\vect{s}_{\uparrow,k})$
in a single factorisation.
\end{remark}

\subsection{Cost analysis}

\begin{itemize}[nosep]
\item \textbf{Per step}: one $(2N\!+\!1)\!\times\!(2N\!+\!1)$ matrix
  exponential ($O(N^3)$ via Pad\'e scaling and
  squaring~\cite{higham2005,moler2003}), two $N\!\times\!N$ solves
  ($O(N^3)$ each), and $O(N^2)$ matrix--matrix and matrix--vector
  products.  Total: $O(N^3)$ per step.
\item \textbf{Overall}: $O(K\cdot M_F\cdot N^3)$, where $K$ is the
  number of Magnus steps and $M_F$ the number of Fourier modes.
\item \textbf{Memory}: $O(N^2)$---only the current accumulated
  operators are stored.  No history is needed.
\end{itemize}

\subsection{Source rescaling}

Following section~1.4 of the Documentation, all boundary sources
and beam intensity are divided by
$\rho=\max(I_0,\,\max(b_{\text{pos}}),\,\max(b_{\text{neg}}))$
before entering the solver, and the output intensities are rescaled
by $\rho$ upon exit.  This prevents numerical issues when source
magnitudes span many orders of magnitude.

\subsection{Fourier mode loop}

The outer loop over Fourier modes $m=0,\dots,M_F-1$ proceeds as
in sections~3.7--3.8 of the Documentation.  Each mode $m$ has its
own coefficient matrix $\mat{A}^m(\tau)$ and source vector
$\vect{S}^m(\tau)$ (via the mode-$m$ phase-function kernel
$\mathcal{D}^m$), and produces its own accumulated scattering
operators.  After the boundary-condition solve (\Cref{sec:bc_solver}),
the full intensity at $\tau=0$ is reconstructed by
\begin{equation}\label{eq:fourier_recon}
  u(\tau\!=\!0,\,\mu_i,\,\phi)
  = \sum_{m=0}^{M_F-1} u^m_i(0)\,\cos\bigl(m(\phi_0-\phi)\bigr).
\end{equation}


%======================================================================
\section{Boundary-Condition Solver}\label{sec:bc_solver}
%======================================================================

After the accumulation loop, the six operators encode the full slab's
input--output relation.  The boundary conditions
(cf.\ section~3.7 of the Documentation) are:
\begin{itemize}[nosep]
\item \emph{Top}: $\vect{u}^-(0)=\vect{b}_{\text{neg}}$
  (prescribed downwelling diffuse intensity at the top of atmosphere).
\item \emph{Bottom}: $\vect{u}^+(\tbot)=\mat{R}_{\text{surf}}\,
  \vect{u}^-(\tbot)+\vect{b}_{\text{pos}}^{\text{eff}}$
  (surface reflection via BDRF plus thermal/isotropic emission).
\end{itemize}

\subsection{BDRF surface matrix}

For Fourier mode $m$, the surface reflection matrix is
\begin{equation}\label{eq:R_surf}
  [\mat{R}_{\text{surf}}]_{ij}
  = (1+\delta_{m0})\,q^m(\mu_i,\mu_j)\,w_j\mu_j,
\end{equation}
where $q^m(\mu,\mu')$ is the $m$-th Fourier mode of the BDRF.  The
factor $(1+\delta_{m0})$ accounts for the $\frac{1}{2}$ in the cosine
series normalisation (cf.\ section~3.7.2 of the Documentation).
If no BDRF is specified, $\mat{R}_{\text{surf}}=\mat{0}$ (black
surface).

The effective bottom source includes the direct-beam surface
reflection:
\begin{equation}\label{eq:bpos_eff}
  \vect{b}_{\text{pos}}^{\text{eff}}
  = \vect{b}_{\text{pos}}
  + \mu_0\,\frac{I_0}{4\pi}\cdot 4\cdot
    q^m(\mu_i,\mu_0)\,e^{-\tbot/\mu_0}\cdot\vect{1},
\end{equation}
where the second term represents the collimated beam reflecting
off the surface (the factor of~4 absorbs the $(2-\delta_{m0})\cdot 2$
from the phase-function convention and Fourier normalisation).

\subsection{The $N\!\times\!N$ boundary system}

Substituting the surface BC into~\eqref{eq:interaction_down}:
\begin{align}
  \vect{u}^+(\tbot)
  &= \mat{R}_{\text{surf}}\,\vect{u}^-(\tbot)
   + \vect{b}_{\text{pos}}^{\text{eff}} \nonumber\\
  &= \mat{R}_{\text{surf}}\,
     \bigl(\mat{T}_\downarrow\,\vect{b}_{\text{neg}}
     + \mat{R}_\downarrow\,\vect{u}^+(\tbot)
     + \vect{s}_\downarrow\bigr)
   + \vect{b}_{\text{pos}}^{\text{eff}}.
\end{align}
Rearranging for $\vect{u}^+(\tbot)$:
\begin{equation}\label{eq:bc_system}
  \bigl(\mat{I}_N - \mat{R}_{\text{surf}}\,\mat{R}_\downarrow\bigr)\,
  \vect{u}^+(\tbot)
  = \mat{R}_{\text{surf}}\,
    \bigl(\mat{T}_\downarrow\,\vect{b}_{\text{neg}}
    + \vect{s}_\downarrow\bigr)
  + \vect{b}_{\text{pos}}^{\text{eff}}.
\end{equation}
This is an $N\!\times\!N$ linear system---\emph{half} the size
of the $2N\!\times\!2N$ Stamnes--Conklin boundary system in the
standard eigendecomposition approach.  It is well-conditioned because
$\|\mat{R}_{\text{surf}}\|\cdot\|\mat{R}_\downarrow\|<1$ for
$\omega<1$.

\subsection{Recovery of the ToA intensity}

Having solved~\eqref{eq:bc_system} for $\vect{u}^+(\tbot)$,
the upward intensity at the top of atmosphere is recovered by
direct substitution into~\eqref{eq:interaction_up}:
\begin{equation}\label{eq:recovery}
  \vect{u}^+(0)
  = \mat{R}_\uparrow\,\vect{b}_{\text{neg}}
  + \mat{T}_\uparrow\,\vect{u}^+(\tbot)
  + \vect{s}_\uparrow.
\end{equation}
The downward intensity at the top is simply
$\vect{u}^-(0)=\vect{b}_{\text{neg}}$ (the prescribed boundary
condition).

\subsection{Flux computation}

The upward diffuse flux at $\tau=0$ is
\begin{equation}
  F_{\text{up}}(0)
  = 2\pi\sum_{i=1}^{N} w_i\,\mu_i\,u^{m=0}_i(0),
\end{equation}
using only the $m\!=\!0$ Fourier mode (higher modes integrate to zero
over azimuth; cf.\ section~3.8 of the Documentation).


%======================================================================
\section{Convergence and Numerical Properties}\label{sec:convergence}
%======================================================================

\subsection{$O(h^2)$ convergence}

The first-order Magnus method with midpoint quadrature has global error
$O(h^2)$, where $h=\tbot/K$.  This means doubling $K$ should reduce
the error by a factor of~4.  \Cref{tab:convergence} confirms this
experimentally for representative test cases spanning thin to very
thick atmospheres.

\begin{table}[h]\centering
\caption{Convergence of the star-product Magnus integrator.  ``Ratio''
is $\text{err}(K)/\text{err}(2K)$; the theoretical value is 4.0 for
$O(h^2)$.}\label{tab:convergence}
\begin{tabular}{llccc}
\toprule
Test & Parameters
  & $K\!=\!400$ & $K\!=\!800$ & Ratio \\
\midrule
2c & $\tbot\!=\!5$, $\omega\!=\!0.5$, Rayleigh
  & $9.2\times 10^{-4}$ & $2.3\times 10^{-4}$ & 4.00 \\
1d & $\tbot\!=\!32$, $\omega\!=\!0.2$, isotropic
  & $4.0\times 10^{-2}$ & $1.0\times 10^{-2}$ & 3.89 \\
1f & $\tbot\!=\!32$, $\omega\!=\!0.99$, isotropic
  & $2.7\times 10^{-2}$ & $7.0\times 10^{-3}$ & 3.94 \\
\bottomrule
\end{tabular}
\end{table}

\subsection{Step-size requirements}

Two constraints determine the minimum number of steps $K$:
\begin{enumerate}[nosep]
\item \textbf{Magnus convergence radius}: $h\cdot\|\mat{A}\|_2 < \pi$,
  i.e.\ $K > \lambda_{\max}\tbot/\pi$.
\item \textbf{Per-step conditioning}: $\cond(\mat{a})$ should remain
  modest; empirically $h\cdot\lambda_{\max}\lesssim 1$--$2$ suffices,
  i.e.\ $K \gtrsim \lambda_{\max}\tbot/2$.
\end{enumerate}
The second constraint is typically the binding one.  For $\tbot=32$ and
$\lambda_{\max}\approx 14$: $K\gtrsim 224$.  In the test suite,
$K=2500$ is used for $\tbot=32$ to achieve $<0.5\%$ relative error.

\subsection{Boundedness verification}

Across all 46 test cases in the test suite (spanning $\tbot$ from
0.03125 to 32, $\omega$ from $10^{-10}$ to $1\!-\!10^{-6}$, isotropic
to HG $g\!=\!0.85$, with and without BDRF and boundary sources), the
following invariants hold at every step of the accumulation:
\begin{itemize}[nosep]
\item $\cond(\mat{a}_k) \leq 1.30$
  (block of each Magnus step),
\item $\cond(\mat{I}-\mat{r}_\uparrow\mat{R}_\downarrow) \leq 1.06$
  (resolvent in star product),
\item $\|\mat{R}_\uparrow\|_2 \leq 0.94$
  (reflection always sub-unitary for $\omega<1$),
\item $\|\mat{T}\|_2 \leq 1.0$
  (transmission bounded).
\end{itemize}
No numerical instability has been observed for any optical depth or
albedo combination.

\subsection{Comparison with pydisort eigendecomposition}

For constant-$\omega$ / constant-phase-function problems (test suites
1--5, 8, 11), the pydisort eigendecomposition solver provides an exact
single-layer reference.  The Magnus solver matches this reference
to within the expected $O(h^2)$ discretisation error:
\begin{itemize}[nosep]
\item Thin atmospheres ($\tbot\leq 2$): relative errors $< 5\times 10^{-3}$
  with $K=100$--$200$.
\item Thick atmospheres ($\tbot=5$--$32$): relative errors $< 5\times 10^{-3}$
  with $K=200$--$2500$.
\end{itemize}

For $\tau$-varying problems (test suites 6--7, 9--10), the Magnus
solver serves as the reference, and the multi-layer pydisort
approximation (piecewise-constant midpoint rule with $N_L$ layers)
converges toward the Magnus solution as $N_L\to\infty$, confirming
consistency between the two approaches.


%======================================================================
\section{Comparison with Related Methods}\label{sec:comparison}
%======================================================================

The star-product Magnus integrator is not the only algorithm for
vertically inhomogeneous atmospheres.  This section compares it with
two established approaches: the piecewise-homogeneous eigendecomposition
used by SCIATRAN and standard DISORT variants, and the adding-doubling
method.

\subsection{First-order Magnus and multi-layer DISORT}\label{sec:magnus_equiv}

A key observation is that the first-order Magnus integrator is
\emph{mathematically equivalent} to a multi-layer discrete-ordinate
solver with midpoint evaluation of optical properties.

In the standard multi-layer DISORT approach~\cite{stamnes1988}, the
atmosphere is divided into $K$ homogeneous layers, each with constant
optical properties evaluated at some representative depth.  Within layer
$k$, the coefficient matrix $\mat{A}_k$ is constant, and the solution
is $\exp(h\mat{A}_k)$.  When the representative depth is the
midpoint $\tau_{k+1/2}$, this is \emph{identical} to the first-order
Magnus step~\eqref{eq:magnus_step}:
\begin{equation}\label{eq:equiv}
  \underbrace{\exp\bigl(h\,\mat{A}(\tau_{k+1/2})\bigr)}_{\text{first-order Magnus}}
  \;=\;
  \underbrace{\exp\bigl(h\,\mat{A}_k\bigr)\big|_{\mat{A}_k\,=\,\mat{A}(\tau_{k+1/2})}}_{\text{multi-layer DISORT, midpoint rule}}.
\end{equation}
Both yield $O(h^2)$ global convergence.  The difference is purely
\emph{implementational}: DISORT computes $\exp(h\mat{A}_k)$ via
eigendecomposition and the Stamnes--Conklin substitution, while the
Magnus integrator uses \texttt{scipy.linalg.expm} (Pad\'e
scaling-and-squaring).  For the $\tau$-independent case, the
eigendecomposition is preferable because a \emph{single} decomposition
yields the exact solution; the Magnus approach should not be used for
such problems.  For $\tau$-varying problems, both implementations yield
the same mathematical approximation at first order.

\subsection{SCIATRAN's approach}\label{sec:sciatran}

The SCIATRAN radiative transfer model~\cite{rozanov2014} is among the
most widely used operational solvers for vertically inhomogeneous
atmospheres in the UV--visible--SWIR range.  Its plane-parallel
discrete-ordinate solver employs:
\begin{enumerate}[nosep]
\item \textbf{Profile discretisation}: optical properties (extinction
  coefficient, single-scattering albedo, phase-function moments) are
  specified at a set of altitude grid points.  Between grid points, the
  properties are interpolated---typically by \emph{linear
  interpolation} in optical depth~\cite{rozanov2014,rozanov2005}.
\item \textbf{Piecewise-constant layers}: each interpolation interval is
  treated as a homogeneous layer with properties evaluated at the
  midpoint (or at the bounding grid point, depending on the
  configuration).
\item \textbf{Eigendecomposition}: within each homogeneous layer, the
  coefficient matrix $\mat{A}$ is diagonalised, yielding the eigenpairs
  $(\lambda_j, \vect{g}_j)$.  The solution within the layer is a linear
  combination of exponentials $e^{\pm\lambda_j\tau}$, with coefficients
  determined by the inter-layer boundary conditions.
\item \textbf{Boundary coupling}: the layers are coupled via a
  banded linear system (the Stamnes--Conklin
  approach~\cite{stamnes1984}) or via matrix-operator/adding
  coupling.
\end{enumerate}

The profile representation by linear splines gives $O(h^2)$
approximation of the underlying continuous profile, which is
well-matched to the $O(h^2)$ error of the piecewise-constant solver.
This is a sound design: a higher-order profile interpolation (e.g.\
cubic splines) would be \emph{wasted} if the solver still treats each
layer as homogeneous, since the solver error would dominate.

\begin{remark}[CDIPI is unrelated]
SCIATRAN's ``CDIPI'' solver~\cite{rozanov2005} (Combined
Differential-Integral approach involving Picard Iterative approximation)
is designed for \emph{spherical geometry}, not for improving the
discrete-ordinate solver in plane-parallel geometry.  It should not be
confused with the matrix-exponential formalism discussed in
\Cref{sec:matexp_rt}.
\end{remark}

\subsection{The adding-doubling method}\label{sec:adding_doubling}

The adding-doubling method~\cite{hansen1974,vandehulst1980,
evans1991} computes the reflection and transmission of a slab by:
\begin{enumerate}[nosep]
\item \textbf{Initialisation}: compute $\mat{R}$ and $\mat{T}$ for a
  very thin slab (small $\Delta\tau$) using a local method---either
  first-principles single-scattering, eigendecomposition, or a
  truncated Peano--Baker (Neumann) series.
\item \textbf{Doubling}: for a homogeneous slab, combine the thin-slab
  operators with themselves via the interaction principle (star product)
  to double the optical depth.  After $s$ doublings, the slab has
  $\Delta\tau\cdot 2^s$ optical thickness.
\item \textbf{Adding}: for an inhomogeneous atmosphere, combine
  \emph{different} thin slabs via the same interaction principle.
\end{enumerate}

The star product used in step~3 is \emph{identical} to the Redheffer
star product of \Cref{sec:star_product}.  The difference from our
Magnus approach is in the initialisation (step~1): adding-doubling
typically uses eigendecomposition or first-order scattering, while we
use the Magnus matrix exponential.

\begin{remark}[Sub-step doubling for large $N$]
For large numbers of streams ($N \gg 4$), the Magnus step size is
constrained by $h \cdot \lambda_{\max} \lesssim 1$, where
$\lambda_{\max} \sim 1/\mu_{\min} \sim O(N^2)$.  An alternative is
\emph{sub-step doubling}: compute the Magnus exponential for a very thin
sub-step where $h'\cdot\lambda_{\max}\lesssim 1$, then self-double via
the star product $s=\lceil\log_2(h/h')\rceil$ times.  This is
precisely the doubling step of adding-doubling, with the Magnus
exponential serving as the initialisation.  It decouples the accuracy
step size (set by the rate of change of $\mat{A}(\tau)$) from the
conditioning step size (set by $\lambda_{\max}$).
\end{remark}

\subsection{The matrix-exponential formalism in radiative transfer}
\label{sec:matexp_rt}

The use of matrix exponentials as a unifying framework for radiative
transfer was developed systematically by Doicu and
Trautmann~\cite{doicu2009}.  Efremenko et
al.~\cite{efremenko2017} proved the formal equivalence of three
approaches under the same optical-property discretisation:
\begin{enumerate}[nosep]
\item the \textbf{discrete-ordinate method} (eigendecomposition of
  $\mat{A}$),
\item the \textbf{matrix-operator method} (adding-doubling with
  $\mat{R}$, $\mat{T}$ operators), and
\item the \textbf{matrix Riccati equation}~\eqref{eq:riccati}
  (invariant imbedding).
\end{enumerate}
All three are different \emph{implementations} of $\exp(h\mat{A})$:
eigendecomposition computes it via diagonalisation, the matrix-operator
method via repeated squaring of thin-slab operators, and the Riccati
equation via the associated M\"obius
transformation~\eqref{eq:mobius}~\cite{efremenko2017}.

The Magnus expansion provides a fourth implementation: direct
Pad\'e approximation of $\exp(\mat{\Omega})$.  Its distinctive
advantage is not at first order (where all methods are equivalent) but
at \emph{higher orders}, where the commutator structure of the Magnus
expansion has no analog in piecewise-constant frameworks
(\Cref{sec:higher_order}).

\subsection{Computational comparison}\label{sec:comp_cost}

\Cref{tab:comparison} summarises the key properties of the three
approaches for a $\tau$-varying atmosphere with $2N$ streams and $K$
steps/layers.

\begin{table}[h]\centering
\caption{Comparison of methods for vertically inhomogeneous atmospheres.}
\label{tab:comparison}
\begin{tabular}{lccc}
\toprule
& \textbf{Multi-layer DISORT} & \textbf{Adding-doubling} & \textbf{Magnus} \\
\midrule
Per-step init.
  & eigendecomp.\ $O(N^3)$
  & eigendecomp.\ $O(N^3)$
  & \texttt{expm}\ $O(N^3)$ \\
Step combination
  & banded solve
  & star product
  & star product \\
Overall cost
  & $O(K N^3)$
  & $O(K N^3)$
  & $O(K N^3)$ \\
Large-$N$ constraint
  & none\textsuperscript{a}
  & $\Delta\tau\!\cdot\!\lambda_{\max}\!\lesssim\!1$
  & $h\!\cdot\!\lambda_{\max}\!\lesssim\!1$\textsuperscript{b} \\
Higher-order ext.
  & Richardson extrap.
  & not standard
  & 4th-order Magnus \\
Profile repr.
  & piecewise const.
  & piecewise const.
  & continuous \\
Adaptive error est.
  & none standard
  & none standard
  & commutator norm \\
\bottomrule
\end{tabular}
\par\smallskip
{\footnotesize\textsuperscript{a}Eigendecomposition is exact for
constant $\mat{A}$; step size limited only by profile
resolution.\quad\textsuperscript{b}Remedied by sub-step doubling
(\Cref{sec:adding_doubling}).}
\end{table}

The Peano--Baker series (the Volterra--Neumann series of nested
integrals) provides yet another representation of the
solution.  However, it converges more slowly than the Magnus expansion
at the same truncation order: the Peano--Baker series at order $p$
is equivalent to the Taylor expansion of $\exp(\mat{\Omega})$, while
the Magnus expansion at order $p$ preserves the Lie-group structure
of the solution manifold (the group of invertible
matrices)~\cite{iserles2000,blanes2009}.  In practice, this means the
Magnus expansion gives better accuracy per function evaluation.


%======================================================================
\section{Higher-Order Magnus and Solver-Function Smoothness}
\label{sec:higher_order}
%======================================================================

\subsection{The fourth-order Magnus integrator}\label{sec:magnus4}

The first-order Magnus method uses a single midpoint evaluation of
$\mat{A}(\tau)$ per step, yielding
$\mat{\Omega}\approx h\mat{A}(\tau_{k+1/2})$.
The \emph{fourth-order} Magnus method~\cite{blanes2009,iserles2000}
uses two Gauss--Legendre quadrature points within each step:
\begin{equation}\label{eq:GL_points}
  \tau_1 = \tau_k + \bigl(\tfrac{1}{2}-\tfrac{\sqrt{3}}{6}\bigr)h,
  \qquad
  \tau_2 = \tau_k + \bigl(\tfrac{1}{2}+\tfrac{\sqrt{3}}{6}\bigr)h,
\end{equation}
and approximates the Magnus exponent as
\begin{equation}\label{eq:magnus4}
  \mat{\Omega} \approx \frac{h}{2}\bigl(\mat{A}_1+\mat{A}_2\bigr)
  + \frac{\sqrt{3}\,h^2}{12}\bigl[\mat{A}_2,\,\mat{A}_1\bigr],
\end{equation}
where $\mat{A}_j=\mat{A}(\tau_j)$ and $[\mat{A}_2,\mat{A}_1]=
\mat{A}_2\mat{A}_1-\mat{A}_1\mat{A}_2$ is the matrix commutator.

The first term, $(h/2)(\mat{A}_1+\mat{A}_2)$, is a two-point
Gauss--Legendre quadrature of $\int\mat{A}\dd\tau$ (which alone would
give a fourth-order Runge--Kutta-type method for a scalar or commuting
$\mat{A}$).  The second term, the commutator correction, captures the
\emph{non-commutativity} of scattering at different depths---i.e.\ the
fact that the order in which radiation encounters different scattering
regimes matters.  The global convergence is $O(h^4)$.

\begin{remark}[No piecewise-constant analog]
The commutator correction~\eqref{eq:magnus4} has no counterpart in
piecewise-constant frameworks (multi-layer DISORT, standard
adding-doubling).  Within a homogeneous layer, $\mat{A}$ is constant,
so $[\mat{A},\mat{A}]=\mat{0}$---the commutator vanishes identically.
Between layers, the inter-layer coupling (boundary conditions or star
product) handles the transition, but it does not capture the continuous
non-commutativity that the Magnus commutator encodes within each step.
To achieve $O(h^4)$ convergence with a piecewise-constant method, one
must take $K\propto h^{-1}\propto\varepsilon^{-1/2}$ layers (for
tolerance $\varepsilon$), whereas the fourth-order Magnus achieves
$K\propto\varepsilon^{-1/4}$---a substantial reduction in the number
of steps for the same accuracy.
\end{remark}

\subsection{Solver-function smoothness matching}\label{sec:smoothness}

A $p$-th order numerical integrator can only achieve its full
convergence rate if the integrand is sufficiently smooth.
Specifically, an $O(h^p)$ Magnus integrator applied to
$\vect{I}'=\mat{A}(\tau)\vect{I}+\vect{S}(\tau)$ requires
$\mat{A}\in C^{p-1}$ to achieve its full convergence rate.  If
$\mat{A}(\tau)$ is only $C^{q-1}$ with $q<p$, the method saturates at
$O(h^q)$ because the higher derivatives appearing in the truncation
error are unbounded or
discontinuous~\cite{blanes2009,hairer2006}.  In particular, the
fourth-order method ($p\!=\!4$) requires $\mat{A}\in C^3$---i.e.\ a
locally cubic profile---since the two-point Gauss--Legendre quadrature
integrates cubic polynomials exactly, and the leading error involves the
fourth derivative of $\mat{A}$.

This creates a \emph{matching principle}: the order of the numerical
solver should be matched to the smoothness of the optical-property
profile representation:
\begin{itemize}[nosep]
\item \textbf{Linear spline profile ($C^0$) + first-order Magnus}:
  $O(h^2)$ convergence.  The profile has slope discontinuities at
  grid points, limiting the useful Taylor expansion of $\mat{A}(\tau)$
  to the constant and linear terms.  A first-order Magnus (or
  equivalently, multi-layer DISORT with midpoint evaluation) is
  perfectly matched.

\item \textbf{Cubic spline profile ($C^2$) + fourth-order Magnus}:
  $O(h^3)$ convergence.  The cubic spline has continuous second
  derivatives but discontinuous third derivatives at knots, so the
  method does not quite achieve its full $O(h^4)$ rate.
  For the full $O(h^4)$, a $C^3$ representation (e.g.\ quintic Hermite
  splines) would be needed.

\item \textbf{$C^3$ or smoother profile + fourth-order Magnus}:
  $O(h^4)$ convergence.  This is the natural pairing for retrieval
  problems where the physical profile is smooth (e.g.\ adiabatic
  clouds, where $\omega(\tau)$ and $g(\tau)$ are $C^\infty$).

\item \textbf{Linear spline + fourth-order Magnus}: the solver achieves
  only $O(h^2)$, because the commutator correction amplifies the slope
  discontinuities at the spline knots, producing no net accuracy gain
  over first-order.  The additional cost of evaluating
  $[\mat{A}_2,\mat{A}_1]$ is wasted.

\item \textbf{Cubic spline + first-order Magnus}: the solver achieves
  $O(h^2)$ despite the profile being $C^2$.  The smooth profile
  information is not exploited; the smooth data justifies a
  higher-order solver that is not being used.
\end{itemize}

This matching principle explains the design of existing solvers:
SCIATRAN's use of linear spline profiles with a piecewise-constant
eigendecomposition solver is well-matched at $O(h^2)$.  Neither
component limits the other.

\subsection{Application to adiabatic cloud profiles}\label{sec:adiabatic}

In an adiabatic liquid water cloud, the liquid water content increases
linearly with height from cloud base~\cite{brenguier2000}, and the
effective radius $r_e$ increases monotonically from cloud top to cloud
base (in optical depth coordinates).  The mapping $r_e\to(\omega,g)$
through Mie scattering theory~\cite{hansen1974} is smooth (analytic,
in fact), so the resulting profiles $\omega(\tau)$ and $g(\tau)$ inherit
the smoothness of $r_e(\tau)$.

For a purely adiabatic cloud with a smooth $r_e(\tau)$ profile:
\begin{itemize}[nosep]
\item $\omega(\tau)$ and $g_\ell(\tau)$ are $C^\infty$ within the cloud.
\item The coefficient matrix $\mat{A}(\tau)$ inherits this smoothness.
\item A fourth-order Magnus integrator achieves its full $O(h^4)$
  convergence, requiring far fewer steps than a first-order method for
  the same accuracy.
\end{itemize}

This is the primary intended use case for higher-order Magnus: cloud
retrieval problems where the microphysical profile is smooth and the
forward model must be evaluated repeatedly (e.g.\ in an iterative
inversion), making the step-count reduction from $O(\varepsilon^{-1/2})$
to $O(\varepsilon^{-1/4})$ operationally significant.

\subsection{Profiles with precipitation perturbations}\label{sec:precip}

When precipitation develops within a cloud (drizzle or rain), large
drops from higher altitudes fall through lower cloud layers, creating a
bimodal drop size distribution.  The effective radius and scattering
properties change abruptly at the precipitation onset
level~\cite{pruppacher1997}, introducing:
\begin{itemize}[nosep]
\item A discontinuity or near-discontinuity in $dr_e/d\tau$
  (the profile is continuous but its derivative is not).
\item A corresponding $C^0$ or worse discontinuity in $\mat{A}(\tau)$.
\end{itemize}

At such boundaries, the fourth-order Magnus commutator correction
encounters a large or undefined $\mat{A}'(\tau)$, and the convergence
degrades to $O(h^2)$ or worse unless the step boundaries are aligned
with the discontinuity.

The practical strategy for mixed profiles is:
\begin{enumerate}[nosep]
\item \textbf{Detect discontinuities} in the prescribed $r_e(\tau)$
  profile (locations where the profile or its derivative changes
  abruptly).
\item \textbf{Align step boundaries} with these discontinuities---i.e.\
  place a star-product interface at each $C^0$ boundary.
\item \textbf{Use higher-order Magnus within smooth segments} and
  first-order across discontinuities.
\end{enumerate}

This mixed-order strategy is a natural extension of the solver and
requires no fundamental changes to the algorithm: the star product
combines arbitrary slabs regardless of how each slab's operators were
computed.


%======================================================================
\section{Toward Adaptive Step Control}\label{sec:adaptivity}
%======================================================================

\subsection{Adaptivity in existing methods}\label{sec:adapt_existing}

Neither SCIATRAN nor standard adding-doubling implementations include
adaptive layer/step placement:
\begin{itemize}[nosep]
\item \textbf{SCIATRAN}~\cite{rozanov2014}: the altitude grid is
  user-specified.  While the user can choose a finer grid in regions of
  rapid variation (e.g.\ near cloud boundaries), the solver itself does
  not adjust the grid automatically.  No built-in local error estimator
  is available.
\item \textbf{Adding-doubling}~\cite{hansen1974}: the standard method
  uses uniform doubling (powers of~2 in optical depth).  For
  inhomogeneous atmospheres, non-uniform adding is possible in
  principle, but there is no standard local error indicator to guide
  step placement.
\end{itemize}

Neither method is fundamentally \emph{prevented} from being
adaptive---the mathematics admits variable step sizes.  The obstacle is
the lack of a cheap, reliable local error estimator.

\subsection{The commutator error indicator}\label{sec:commutator}

The Magnus expansion provides a natural and essentially free error
indicator: the commutator
\begin{equation}\label{eq:commutator_indicator}
  \eta_k = h_k^2\,\bigl\|\bigl[\mat{A}(\tau_k),\;\mat{A}(\tau_k+h_k)\bigr]\bigr\|_F,
\end{equation}
where $\|\cdot\|_F$ is the Frobenius norm.

\textbf{Interpretation.}  The commutator measures how much $\mat{A}$
changes \emph{direction} (not just magnitude) across a step:
\begin{itemize}[nosep]
\item \textbf{Exactly zero} when $\mat{A}$ is constant---i.e.\ in a
  homogeneous atmosphere, or within a homogeneous layer.  In such
  regions, arbitrarily large steps are safe (up to the Magnus
  convergence radius).
\item \textbf{Small} when $\mat{A}$ varies slowly and smoothly---e.g.\
  in the interior of an adiabatic cloud, where $r_e(\tau)$ changes
  gradually.
\item \textbf{Large} at cloud boundaries, precipitation onset levels,
  or any region where the scattering regime changes rapidly.
\end{itemize}

For the fourth-order Magnus~\eqref{eq:magnus4}, the commutator
$[\mat{A}_2,\mat{A}_1]$ is \emph{already computed} as part of the
step.  Its norm therefore provides an error indicator at
\emph{zero additional cost}.  Even for the first-order Magnus, evaluating
$[\mat{A}(\tau_k),\mat{A}(\tau_k+h_k)]$ requires only two matrix
multiplications ($O(N^3)$), which is negligible compared to the
\texttt{expm} cost.

\subsection{Embedded-pair error estimation}\label{sec:embedded_pair}

The most rigorous approach to adaptive step control uses an
\emph{embedded pair}: two methods of different order applied to the same
step, whose difference provides a direct local truncation error
estimate.  For Magnus integration, a natural embedded pair combines the
first-order (midpoint) and fourth-order (Gauss--Legendre) methods.

Given a step $[\tau_k,\,\tau_k+h_k]$, the first-order Magnus exponent
is
\begin{equation}\label{eq:magnus1_embed}
  \mat{\Omega}_2 = h_k\,\mat{A}\bigl(\tau_k+\tfrac{h_k}{2}\bigr),
\end{equation}
while the fourth-order exponent~\eqref{eq:magnus4} is
\begin{equation}\label{eq:magnus4_embed}
  \mat{\Omega}_4 = \frac{h_k}{2}(\mat{A}_1+\mat{A}_2)
  + \frac{\sqrt{3}\,h_k^2}{12}[\mat{A}_2,\mat{A}_1],
\end{equation}
where $\mat{A}_1,\mat{A}_2$ are evaluated at the two Gauss--Legendre
points~\eqref{eq:GL_points}.  Both exponents are computed from the same
step, and the local truncation error is estimated by
\begin{equation}\label{eq:embed_error}
  \varepsilon_k = \bigl\|\operatorname{expm}(\mat{\Omega}_4)
  - \operatorname{expm}(\mat{\Omega}_2)\bigr\|.
\end{equation}

\textbf{Cost.}  The fourth-order method requires two $\mat{A}$
evaluations (at the Gauss--Legendre points) plus one commutator (two
$N\!\times\!N$ multiplications).  The first-order method adds one
$\mat{A}$ evaluation (at the midpoint) and one \texttt{expm} call.  The
total per-step overhead is approximately $2\times$ the cost of a single
first-order step.  However, since the fourth-order result is used to
\emph{advance} the solution (not merely to estimate the error), each
accepted step achieves $O(h^4)$ accuracy, and far fewer steps are needed
in smooth regions.

\textbf{Relation to the commutator indicator.}  When the dominant
source of error is the commutator correction, the embedded error
$\varepsilon_k$ is well-approximated by
$\|\operatorname{expm}(\mat{\Omega}_4)
-\operatorname{expm}(\mat{\Omega}_2)\|\approx C\,\eta_k$, where
$\eta_k$ is the commutator indicator~\eqref{eq:commutator_indicator}.
In such regions, $\eta_k$ is a reliable and cheaper proxy for
$\varepsilon_k$.  The full embedded error is needed only where the
commutator-based estimate is unreliable (e.g.\ near discontinuities in
$\mat{A}$, where higher-order terms dominate).

\subsection{Adaptive Magnus step control}\label{sec:adapt_magnus}

Based on the embedded-pair error estimator, we outline an adaptive
strategy for the Magnus integrator (not yet implemented):

\begin{enumerate}
\item \textbf{Step-size ceiling}: regardless of the accuracy criterion,
  the step size must satisfy the \emph{stability} constraint from
  \Cref{sec:convergence}:
  \begin{equation}\label{eq:h_max}
    h_k \;\leq\; h_{\max} \;=\; \frac{c}{\lambda_{\max}},
    \qquad c \approx 1\text{--}2,
  \end{equation}
  where $\lambda_{\max}\approx 1/\mu_{\min}$ is the spectral radius of
  $\mat{A}$ (\Cref{sec:convergence}).  This ceiling is a hard constraint
  that cannot be relaxed by any accuracy criterion---even in a perfectly
  homogeneous atmosphere ($\eta_k=0$), the step size may not exceed
  $h_{\max}$.

\item \textbf{Step acceptance}: at each step, compute the embedded
  error~\eqref{eq:embed_error}.
  \begin{itemize}[nosep]
  \item If $\varepsilon_k < \varepsilon_{\text{tol}}$: \textbf{accept}
    the step (using the fourth-order result to advance the solution) and
    attempt to \textbf{grow} the next step to
    $\min(\alpha\,h_k,\,h_{\max})$ with $\alpha\in(1,2]$.
  \item If $\varepsilon_k > \varepsilon_{\text{tol}}$: \textbf{reject}
    the step, halve $h$, and retry.
  \end{itemize}
  Where the commutator indicator $\eta_k$ is a reliable proxy (smooth
  regions), it may be used in place of the full embedded error to avoid
  the second \texttt{expm} call.

\item \textbf{Order selection}: in smooth regions (small
  $\varepsilon_k$), use fourth-order Magnus for maximum efficiency.
  Near detected discontinuities (large $\varepsilon_k$ that does not
  decrease sufficiently under halving), fall back to first-order Magnus,
  which remains $O(h^2)$ for any $\mat{A}\in L^1$.

\item \textbf{Star-product accumulation}: the Redheffer star product is
  agnostic to step size---each step produces slab operators that are
  combined regardless of their width.  Variable step sizes and
  mixed-order steps require no changes to the star product machinery.
\end{enumerate}

The effective step size at each point is
\begin{equation}\label{eq:h_effective}
  h_k = \min\bigl(h_{\text{accuracy}}(\varepsilon_k,
  \varepsilon_{\text{tol}}),\;h_{\max}\bigr).
\end{equation}
In smooth regions (e.g.\ the interior of an adiabatic cloud), the
accuracy criterion allows large steps but the stability ceiling
$h_{\max}$ is binding.  In rapidly varying regions (e.g.\ cloud
boundaries), the accuracy criterion demands refinement well below
$h_{\max}$, and the stability ceiling is inactive.

\subsection{The diffusion-domain stability ceiling}
\label{sec:diffusion_ceiling}

The interplay between accuracy and stability
in~\eqref{eq:h_effective} reveals a fundamental limitation of Magnus
integration in the \emph{diffusion domain}---regions where the
atmosphere is optically thick and highly scattering ($\omega\to 1$).

The coefficient matrix $\mat{A}(\tau)$ has eigenvalues $\pm k_j$ where
$k_{\max}\approx 1/\mu_{\min}$ (ballistic modes) and
$k_{\min}\approx\sqrt{3(1-\omega)(1-\omega g)}$ (diffusion mode).  In
the diffusion domain, the accuracy of the solution is governed by
$k_{\min}$: the physical observable (diffuse reflectance) converges at a
rate proportional to $\exp(-2k_{\min}\tau)$, so steps of size
$h\sim 1/k_{\min}$ suffice for accuracy.  However, the Magnus stability
ceiling~\eqref{eq:h_max} enforces $h\leq c/\lambda_{\max}$ with
$\lambda_{\max}\approx k_{\max}$, which is independent of $\omega$.

The \emph{eigenvalue gap ratio}
\begin{equation}\label{eq:gap_ratio}
  \gamma = \frac{k_{\max}}{k_{\min}}
  \approx \frac{1/\mu_{\min}}{\sqrt{3(1-\omega)(1-\omega g)}}
\end{equation}
measures the severity of this mismatch.  For $\omega=0.999$, $g=0.85$,
$N_{\text{Quad}}=8$: $k_{\max}\approx 14$, $k_{\min}\approx 0.067$,
giving $\gamma\approx 210$.  The Magnus integrator must take at least
$K\geq\lambda_{\max}\tau_{\text{bot}}$ steps, even though only
$K\sim k_{\min}\tau_{\text{bot}}$ steps would suffice for accuracy.  The
adaptive strategy of \Cref{sec:adapt_magnus} cannot help: the stability
ceiling $h_{\max}$ is the binding constraint, and no amount of
accuracy-based step selection can overcome it.

This observation motivates a fundamentally different integration
approach in the diffusion domain, developed in detail in
Report~II~\cite{report2}.

\subsection{Hybrid adaptive strategy}\label{sec:hybrid_adaptive}

The limitations identified in \Cref{sec:diffusion_ceiling} lead
naturally to a \emph{hybrid} strategy that uses the most efficient
method in each regime:
\begin{itemize}[nosep]
\item \textbf{Non-diffusive regions} (beam-dominated, cloud boundaries,
  optically thin layers): adaptive fourth-order Magnus with
  embedded-pair step control (\Cref{sec:adapt_magnus}).  The beam
  source $\propto e^{-\tau/\mu_0}$ varies on scale $\mu_0$, the
  accuracy criterion is binding, and the stability ceiling $h_{\max}$
  is typically inactive.
\item \textbf{Diffusion domain} (deep interior of thick, highly
  scattering slabs): the matrix Riccati equation
  \begin{equation}\label{eq:riccati_preview}
    \frac{d\mat{R}}{d\sigma}
    = \mat{\alpha}\mat{R} + \mat{R}\mat{\alpha}
      + \mat{R}\mat{\beta}\mat{R} + \mat{\beta},
    \qquad \mat{R}(0) = \mat{0},
  \end{equation}
  where $\sigma$ is optical depth measured from the bottom of the
  domain, is solved with a Rosenbrock--Euler method that exploits the
  Sylvester structure of the linearised
  equation~\cite{hairer1996,bartels1972}.  The Jacobian of
  the Riccati right-hand side has the Kronecker-sum form
  $\mat{J}(\mat{E})
  = (\mat{\alpha}+\mat{R}\mat{\beta})\mat{E}
  + \mat{E}(\mat{\alpha}+\mat{\beta}\mat{R})$
  (a Sylvester operator), so each implicit step reduces to a
  Sylvester equation solvable in $O(N^3)$ via Bartels--Stewart.
  The resulting integrator is L-stable: its step size is limited only
  by accuracy ($h\sim 1/k_{\min}$), not by the spectral
  radius~$\lambda_{\max}$.
  An embedded pair (Rosenbrock--Euler vs.\ Richardson extrapolation)
  provides adaptive step-size control analogous to the Magnus embedded
  pair.  Report~II~\cite{report2} develops this approach in detail and
  benchmarks it against the Magnus integrator, demonstrating a
  $\sim\!14\!\times$ reduction in step count for representative cloud
  parameters ($\omega=0.99$, $g=0.85$, $\tau_{\text{sub}}=20$,
  $N_{\text{Quad}}=32$).
\end{itemize}

The two domains are coupled via the Redheffer star product, which
combines arbitrary slab operators regardless of how they were computed.
Each domain independently produces its
$(\mat{R}_{\text{up}},\mat{T}_{\text{up}},\mat{T}_{\text{down}},
\mat{R}_{\text{down}},\vect{s}_{\text{up}},\vect{s}_{\text{down}})$
operators, and the star product assembles the full-atmosphere response.
This is a natural extension of the framework of
\Cref{sec:star_product}: the product does not depend on the internal
method used within each slab.

\textbf{Automatic domain detection.}  The transition between regimes
can be detected automatically during integration: when the adaptive
Magnus step controller repeatedly hits the stability ceiling $h_{\max}$
while the embedded error $\varepsilon_k$ is far below tolerance, this
signals that the step count is stability-limited rather than
accuracy-limited---precisely the signature of the diffusion domain.

\subsection{Expected benefits}\label{sec:adapt_benefits}

For the target application---cloud microphysics retrieval with adiabatic
profiles---the hybrid adaptive strategy provides:
\begin{itemize}[nosep]
\item \textbf{Automatic coarsening} in smooth regions, where
  $\omega(\tau)$ and $g(\tau)$ vary slowly and large steps suffice
  (limited by either the accuracy criterion or the stability ceiling,
  whichever is tighter).
\item \textbf{Automatic refinement} near cloud boundaries and
  precipitation onset levels, where optical properties change rapidly
  and the accuracy criterion demands small steps.
\item \textbf{Diffusion-domain efficiency}: in the deep interior of
  thick clouds where $\omega\to 1$, the Rosenbrock--Sylvester integrator
  breaks through the Magnus stability ceiling, reducing the step count
  by a factor of up to~$\gamma$.  For $\omega=0.999$, $g=0.85$,
  $N_{\text{Quad}}=8$, this is a $\sim\!200\!\times$ reduction in the
  diffusion domain (Report~II, \S8.5).
\item \textbf{Robustness}: the solver adjusts to the profile complexity
  without user intervention, unlike SCIATRAN where the grid must be
  manually tuned.  Even if the diffusion-domain solver is inaccurate,
  the star-product coupling bounds errors at $O(\varepsilon)$ rather
  than $O(e^{k_{\max}\tau}\varepsilon)$---there is no catastrophic
  failure mode.
\item \textbf{Compound savings in retrieval}: since the forward solver
  sits inside an iterative retrieval loop, the step-count reduction
  compounds across evaluations, potentially dominating total computation
  time.
\end{itemize}


%======================================================================
\section{Deferred Features}\label{sec:deferred}
%======================================================================

The following features are intentionally deferred from the current
implementation:
\begin{itemize}[nosep]
\item \textbf{Adaptive step control}: currently the $K$ equidistant
  steps are specified by the user.  \Cref{sec:adaptivity} outlines an
  embedded-pair adaptive strategy for Magnus, and a hybrid scheme
  coupling Magnus with Rosenbrock--Sylvester integration in the
  diffusion domain; implementation is deferred.
\item \textbf{Fourth-order Magnus}: the commutator-corrected
  integrator~\eqref{eq:magnus4} is derived in \Cref{sec:magnus4} but
  not yet implemented; only the first-order method is available.
\item \textbf{Delta-$M$ scaling}~\cite{wiscombe1977}: not applied in
  \texttt{pydisort\_magnus}.  For strongly forward-peaked phase
  functions, the effective $\lambda_{\max}$ is smaller after delta-$M$
  scaling, reducing the required $K$.
\item \textbf{Nakajima--Tanaka corrections}~\cite{NT1988}: the TMS/IMS
  post-hoc intensity corrections (section~3.7.3 of the Documentation)
  are not applied.
\item \textbf{Isotropic internal source}: only the collimated beam is
  handled; thermal emission would enter as an additional source term
  in~\eqref{eq:source}.
\item \textbf{Non-ToA evaluation}: currently only $\tau\!=\!0$
  intensities are returned.  Evaluation at intermediate depths would
  require storing per-step operators and performing a partial
  star-product accumulation.
\end{itemize}


%======================================================================
\section{Summary}\label{sec:summary}
%======================================================================

The star-product Magnus integrator resolves the fundamental numerical
stability problem that prevents na\"ive propagator accumulation from
working in thick atmospheres.  Its key properties are:

\begin{enumerate}[nosep]
\item \textbf{Unconditional stability}: all intermediate quantities
  are $O(1)$, regardless of total optical depth.
  No positive exponents appear anywhere.
\item \textbf{No eigendecomposition}: the algorithm uses only matrix
  exponentials (via \texttt{expm}) and linear solves---no
  diagonalisation of $\mat{A}$ is needed at any point.
\item \textbf{$O(h^2)$ convergence}: the first-order Magnus midpoint
  rule gives clean second-order convergence, confirmed experimentally
  with convergence ratios $\approx 4.0$.
\item \textbf{Native support for $\tau$-varying properties}: the
  coefficient matrix $\mat{A}(\tau)$ and source vector $\vect{S}(\tau)$
  are evaluated at the midpoint of each step, naturally handling
  continuously varying $\omega(\tau)$ and phase function.
\item \textbf{$N\!\times\!N$ boundary system}: the BC solve
  is half the size of the standard Stamnes--Conklin system, and
  naturally handles BDRF, $b_{\text{pos}}$, and $b_{\text{neg}}$.
\item \textbf{Cost}: $O(K\cdot M_F\cdot N^3)$ with $O(N^2)$ memory.
\end{enumerate}

At first order, the Magnus integrator is mathematically equivalent to a
multi-layer discrete-ordinate solver with midpoint evaluation
(\Cref{sec:magnus_equiv}), sharing the same $O(h^2)$ convergence as
SCIATRAN's piecewise-constant approach.  The genuine advantages of the
Magnus framework emerge at higher orders: the fourth-order
method~\eqref{eq:magnus4} captures non-commutativity of scattering at
different depths via the matrix commutator, achieving $O(h^4)$
convergence with no piecewise-constant analog.  This higher-order
extension is most effective when matched with sufficiently smooth
profiles (\Cref{sec:smoothness})---for example, the $C^\infty$
profiles arising from adiabatic cloud microphysics.  The commutator provides a natural
adaptive error indicator (\Cref{sec:commutator}), and an embedded
order-2/4 pair (\Cref{sec:embedded_pair}) gives a rigorous local
truncation error estimate for automatic step-size control---a capability
unavailable in standard DISORT or adding-doubling implementations.
In the diffusion domain, where the Magnus stability ceiling is the
binding constraint, a hybrid strategy coupling Magnus with
Rosenbrock--Sylvester integration (\Cref{sec:hybrid_adaptive}) breaks
through this ceiling, potentially reducing step counts by orders of
magnitude for highly scattering atmospheres.


%======================================================================
\begin{thebibliography}{99}

\bibitem{HP2024doc}
D.J.X.~Ho, ``Pythonic-DISORT: Comprehensive Documentation,''
\texttt{docs/Pythonic-DISORT.ipynb} in the PythonicDISORT repository,
2024.
\url{https://pythonic-disort.readthedocs.io}

\bibitem{HP2024}
D.J.X.~Ho and R.~Pincus, ``Two-Streams Revisited: General Equations,
Exact Coefficients, and Optimized Closures,'' \emph{J.\ Adv.\ Model.\
Earth Syst.}, 16(10):e2024MS004504, 2024.
\url{https://doi.org/10.1029/2024MS004504}

\bibitem{magnus1954}
W.~Magnus, ``On the exponential solution of differential equations for
a linear operator,'' \emph{Comm.\ Pure Appl.\ Math.}, 7:649--673,
1954.

\bibitem{blanes2009}
S.~Blanes, F.~Casas, J.A.~Oteo, and J.~Ros, ``The Magnus expansion
and some of its applications,'' \emph{Physics Reports},
470(5--6):151--238, 2009.

\bibitem{moan2008}
P.C.~Moan and J.~Niesen, ``Convergence of the Magnus series,''
\emph{Found.\ Comput.\ Math.}, 8(3):291--301, 2008.

\bibitem{stamnes1988}
K.~Stamnes, S.-C.~Tsay, W.~Wiscombe, and K.~Jayaweera, ``Numerically
stable algorithm for discrete-ordinate-method radiative transfer in
multiple scattering and emitting layered media,'' \emph{Appl.\ Opt.},
27(12):2502--2509, 1988.

\bibitem{stamnes1984}
K.~Stamnes and P.~Conklin, ``A new multi-layer discrete ordinate
approach to radiative transfer in vertically inhomogeneous
atmospheres,'' \emph{J.\ Quant.\ Spectrosc.\ Radiat.\ Transfer},
31(3):273--282, 1984.

\bibitem{redheffer1962}
R.~Redheffer, ``On the relation of transmission-line theory to
scattering and transfer,'' \emph{J.\ Math.\ Phys.}, 41:1--41, 1962.

\bibitem{grant1969}
I.P.~Grant and G.E.~Hunt, ``Discrete space theory of radiative
transfer.\ I.~Fundamentals,'' \emph{Proc.\ R.\ Soc.\ Lond.\ A},
313(1513):183--197, 1969.

\bibitem{plass1973}
G.N.~Plass, G.W.~Kattawar, and F.E.~Catchings, ``Matrix operator
theory of radiative transfer.\ 1: Rayleigh scattering,''
\emph{Appl.\ Opt.}, 12(2):314--329, 1973.

\bibitem{bellman1960}
R.~Bellman, R.~Kalaba, and M.~Prestrud, \emph{Invariant Imbedding and
Radiative Transfer in Slabs of Finite Thickness}, American Elsevier,
1960.

\bibitem{kreiss1962}
H.-O.~Kreiss, ``\"Uber die Stabilit\"atsdefinition f\"ur
Differenzengleichungen die partielle Differentialgleichungen
approximieren,'' \emph{BIT}, 2:153--181, 1962.

\bibitem{vanloan1978}
C.F.~Van~Loan, ``Computing integrals involving the matrix
exponential,'' \emph{IEEE Trans.\ Automat.\ Control},
23(3):395--404, 1978.

\bibitem{moler2003}
C.~Moler and C.~Van~Loan, ``Nineteen dubious ways to compute the
exponential of a matrix, twenty-five years later,'' \emph{SIAM
Review}, 45(1):3--49, 2003.

\bibitem{higham2005}
N.J.~Higham, ``The scaling and squaring method for the matrix
exponential revisited,'' \emph{SIAM J.\ Matrix Anal.\ Appl.},
26(4):1179--1193, 2005.

\bibitem{wiscombe1977}
W.J.~Wiscombe, ``The Delta--$M$ Method: Rapid Yet Accurate Radiative
Flux Calculations for Strongly Asymmetric Phase Functions,''
\emph{J.\ Atmos.\ Sci.}, 34(9):1408--1422, 1977.

\bibitem{NT1988}
T.~Nakajima and M.~Tanaka, ``Algorithms for radiative intensity
calculations in moderately thick atmospheres using a truncation
approximation,'' \emph{J.\ Quant.\ Spectrosc.\ Radiat.\ Transfer},
40(1):51--69, 1988.

\bibitem{chandrasekhar1960}
S.~Chandrasekhar, \emph{Radiative Transfer}, Dover, 1960.

\bibitem{rozanov2014}
V.V.~Rozanov, A.V.~Rozanov, A.A.~Kokhanovsky, and J.P.~Burrows,
``Radiative transfer through terrestrial atmosphere and ocean: Software
package SCIATRAN,'' \emph{J.\ Quant.\ Spectrosc.\ Radiat.\ Transfer},
133:13--71, 2014.

\bibitem{rozanov2005}
A.~Rozanov, V.~Rozanov, M.~Buchwitz, A.~Kokhanovsky, and
J.P.~Burrows, ``SCIATRAN 2.0---A new radiative transfer model for
geophysical applications in the 175--2400~nm spectral region,''
\emph{Adv.\ Space Res.}, 36(5):1015--1019, 2005.

\bibitem{hansen1974}
J.E.~Hansen and L.D.~Travis, ``Light scattering in planetary
atmospheres,'' \emph{Space Sci.\ Rev.}, 16(4):527--610, 1974.

\bibitem{vandehulst1980}
H.C.~van~de~Hulst, \emph{Multiple Light Scattering: Tables, Formulas,
and Applications}, vols.\ 1--2, Academic Press, 1980.

\bibitem{evans1991}
K.F.~Evans and G.L.~Stephens, ``A new polarized atmospheric radiative
transfer model,'' \emph{J.\ Quant.\ Spectrosc.\ Radiat.\ Transfer},
46(5):413--423, 1991.

\bibitem{doicu2009}
A.~Doicu and T.~Trautmann, ``Discrete-ordinate method with matrix
exponential for a pseudo-spherical atmosphere: Scalar case,''
\emph{J.\ Quant.\ Spectrosc.\ Radiat.\ Transfer},
110(1--2):146--158, 2009.

\bibitem{efremenko2017}
D.S.~Efremenko, V.~Molina~Garc\'\i a, S.~Gimeno~Garc\'\i a, and
A.~Doicu, ``A review of the matrix-exponential formalism in radiative
transfer,'' \emph{J.\ Quant.\ Spectrosc.\ Radiat.\ Transfer},
196:17--45, 2017.

\bibitem{iserles2000}
A.~Iserles, H.Z.~Munthe-Kaas, S.P.~N\o rsett, and A.~Zanna,
``Lie-group methods,'' \emph{Acta Numerica}, 9:215--365, 2000.

\bibitem{hairer2006}
E.~Hairer, C.~Lubich, and G.~Wanner, \emph{Geometric Numerical
Integration}, 2nd ed., Springer, 2006.

\bibitem{brenguier2000}
J.-L.~Brenguier, H.~Pawlowska, L.~Sch\"uller, R.~Preusker,
J.~Fischer, and Y.~Fouquart, ``Radiative properties of boundary layer
clouds: Droplet effective radius versus number concentration,''
\emph{J.\ Atmos.\ Sci.}, 57(6):803--821, 2000.

\bibitem{pruppacher1997}
H.R.~Pruppacher and J.D.~Klett, \emph{Microphysics of Clouds and
Precipitation}, 2nd ed., Kluwer Academic Publishers, 1997.

\bibitem{report2}
D.~Halberstadt, ``Diffusion-Domain Decomposition for the Star-Product
Magnus Integrator,'' companion report (Report~II), 2026.

\bibitem{hairer1996}
E.~Hairer and G.~Wanner, \emph{Solving Ordinary Differential
Equations~II: Stiff and Differential-Algebraic Problems}, 2nd ed.,
Springer, 1996.

\bibitem{bartels1972}
R.H.~Bartels and G.W.~Stewart, ``Solution of the matrix equation
$AX+XB=C$,'' \emph{Comm.\ ACM}, 15(9):820--826, 1972.

\end{thebibliography}

\end{document}
